% Unregistered candidate D.58.  Root audit requested before integration.
\paragraph{Candidate D.58: the full mixed fibre under a precise Selmer hypothesis.}

Use the actual cubic twist and the preferred origin of D.8:
\[
 E_c:\ y^3=cx(1-x),\qquad O_c\text{ the cusp over }0,
 \qquad V_c=E_c\setminus\{\text{the three cusps}\}.
\]
All identifications below use that twist convention; replacing the
convention by its inverse merely relabels the cubic class.  Let
$\mathcal F_c\subset H^1(\mathbb Q,E_c[5])$ be the set of
classes in D.8's fibre which extend to a section of the full
geometric tripod group.  Put
\[
 \mathcal S_5(E_c)=\varprojlim_m\operatorname{Sel}_{5^m}(E_c/\mathbb Q),
 \qquad
 R_5(E_c)=\operatorname{im}\bigl(\mathcal S_5(E_c)
                  \longrightarrow H^1(\mathbb Q,E_c[5])\bigr),
\]
where the transition maps on torsion coefficients are multiplication
by five.  Then
\[
 \kappa_5\bigl(E_c(\mathbb Q)/5E_c(\mathbb Q)\bigr)
       \ \subset\ \mathcal F_c\ \subset\ R_5(E_c)
       \ \subset\ \operatorname{Sel}_5(E_c/\mathbb Q).       \tag{D.58.1}
\]
In particular the right hand filters are finite.  There is an exact
sequence of compact abelian groups
\[
 0\longrightarrow E_c(\mathbb Q)^{\wedge 5}
 \longrightarrow\mathcal S_5(E_c)
 \longrightarrow T_5\operatorname{Sha}(E_c/\mathbb Q)\longrightarrow0.  \tag{D.58.2}
\]
Consequently, if $T_5\operatorname{Sha}(E_c/\mathbb Q)=0$, then
\[
 \mathcal F_c=
 \kappa_5\bigl(E_c(\mathbb Q)/5E_c(\mathbb Q)\bigr).       \tag{D.58.3}
\]
The rational five-torsion group of every $E_c$ is zero.  Thus the
set in (D.58.3) has $5^{\operatorname{rank}E_c(\mathbb Q)}$
elements.  In particular (D.58.3) holds whenever $E_c$ has
analytic rank zero or one, with respectively one or five classes.
All classes on the left are realized by rational ordinary points
or rational cusp packets of the actual twist.  This classifies
the first-five image of the full section set under the displayed
hypothesis.  It does not assert a classification of the full
sections themselves, or a classification for the relative
pro-five quotient.

\begin{proof}
\paragraph{1. The actual twisted cover and its proper Kummer image.}

Represent the fixed cubic cohomology class by the actual cocycle
used in D.8.  The inverse image of its graph in the arithmetic
tripod group is the arithmetic fundamental group of the
corresponding twist $V_c$.  A full section with that cubic
class can be conjugated so that its cubic cocycle equals this
representative; its image then lies in this covering subgroup.
Conversely every section of $V_c$ maps to a full tripod section
in the prescribed cubic fibre.  This is the finite etale
covering correspondence with the actual arithmetic twist,
not an inference from a maximal pro-five quotient.

Compactifying $V_c$ gives the pointed elliptic curve $E_c$.
Its geometric fundamental group is $TE_c$, so composing a
lifted section with the proper pro-five abelian quotient
gives a class
\[
                         z_s\in H^1(\mathbb Q,T_5E_c).
\]
At every finite place $v$, D.56 with the formation of all
finite groups makes the lifted full affine section point
or cusp induced.  Its proper image is therefore the Kummer
class of some $P_v\in E_c(\mathbb Q_v)$.  At the real place,
$H^1(\mathbb R,E_c[5^m])=0$ for every $m$, by averaging over
the group of order two; there is no missing five-primary
local condition.  Reducing $z_s$ modulo $5^m$ at each $m$
therefore gives a compatible class in
$\operatorname{Sel}_{5^m}(E_c/\mathbb Q)$.
Its first reduction is precisely the translation coordinate
of the section in D.8, by the origin and twist comparison
there.  This proves the middle containment in (D.58.1).

For the reverse containment on the left, take any
$P\in E_c(\mathbb Q)$.  If $P\in V_c$, its ordinary point
section gives the Kummer class $\kappa_5(P-O_c)$.  If $P$
is one of the three deleted points, a nonzero rational
tangent gives a cusp section with the same proper Kummer
class.  All three cusps are rational, and $O_c$ is the
origin, so this construction applies to every rational
point of $E_c$.  It proves the first containment in (D.58.1).

\paragraph{2. The inverse limit and the exact Shafarevich--Tate term.}

The finite Kummer sequences and the definition of the Selmer
local conditions give, for each $m\ge1$,
\[
 0\longrightarrow E_c(\mathbb Q)/5^mE_c(\mathbb Q)
 \longrightarrow\operatorname{Sel}_{5^m}(E_c/\mathbb Q)
 \longrightarrow\operatorname{Sha}(E_c/\mathbb Q)[5^m]
 \longrightarrow0.                                  \tag{D.58.4}
\]
The same sequence and finiteness of its middle term are
displayed in Kolyvagin's primary manuscript, Section 2,
equation (5), printed page 6 \cite{DE_Kolyvagin1990}.
The left hand transition maps are the natural quotient
maps and are surjective.  Thus this inverse system is
Mittag--Leffler, and its first derived inverse limit is
zero.  Taking inverse limits of (D.58.4) proves (D.58.2),
including surjectivity on the right.  All finite Selmer
groups carry their finite discrete topology; their inverse
limit is compact.  There is no assertion that the global
groups $H^1(\mathbb Q,E_c[5^m])$ are finite, and no
unjustified omission of a global derived-limit term.

If $T_5\operatorname{Sha}=0$, equation (D.58.2) identifies the inverse
Selmer group with the five-adic completion of the
Mordell--Weil group.  Its reduction modulo five has image
exactly $E_c(\mathbb Q)/5E_c(\mathbb Q)$, since a finitely
generated abelian group surjects onto this finite quotient.
The two containments already proved now give (D.58.3).
Notice that finite five-primary Shafarevich--Tate classes
may occur in $\operatorname{Sel}_5$; they do not survive
as a nonzero compatible Tate-module contribution.  This
is why the inverse Selmer group is the correct filter.

\paragraph{3. The torsion count and the unconditional analytic-rank cases.}

Choose a prime $q\equiv17\pmod{30}$ at which $c$ is a unit.
Such a prime exists by Dirichlet's theorem, and it is a
good reduction prime for the plane cubic $E_c$, since
$q\nmid3c$.  Because $q\equiv2\pmod3$, cubing is a
bijection on $\mathbb F_q$.  For each of the $q$ affine
values of $x$ there is exactly one $y$, and there is one
point at infinity.  Therefore
\[
                         |E_c(\mathbb F_q)|=q+1
                                      \equiv3\pmod5.
\]
Good reduction injects rational five-torsion into this
group, so $E_c(\mathbb Q)[5]=0$.  The structure theorem
for the finitely generated Mordell--Weil group gives the
stated cardinality in (D.58.3).

Every elliptic curve over $\mathbb Q$ is modular by
\cite[Theorem A]{DE_BCDT2001}.  Kolyvagin's primary
1990 manuscript, Theorem 1, printed page 4, states that
a modular elliptic curve of analytic rank at most one
has equal algebraic and analytic ranks and finite
Shafarevich--Tate group \cite{DE_Kolyvagin1990}.
The next page explains the Gross--Zagier and
quadratic-twist nonvanishing inputs that remove the
auxiliary Heegner-point hypothesis of the earlier theorem.
Thus $T_5\operatorname{Sha}(E_c/\mathbb Q)=0$ in these cases, with no
restriction that five be an ordinary or good reduction
prime.  Equations (D.58.2)--(D.58.3) finish the proof.
\end{proof}

\paragraph{Remaining scope of Task D.9.}

For arbitrary $c$ with a possible nonzero
$T_5\operatorname{Sha}(E_c/\mathbb Q)$, this argument supplies the
finite inverse-Selmer image $R_5(E_c)$ but does not decide
which of its non-Mordell--Weil classes lift nonabelianly.
It also leaves the relative pro-five question open even
when (D.58.3) settles the full-group first-five image.
